\documentclass[11pt,a4paper]{article}

\usepackage[utf8]{inputenc}
\usepackage[T1]{fontenc}
\usepackage[french]{babel}
\usepackage{amsmath,amssymb,amsfonts,amsthm}
\usepackage{geometry}
\geometry{a4paper, top=2.5cm, bottom=2.5cm, left=2.3cm, right=2.3cm, centering}

% Centrage vertical automatique du contenu de chaque page
\makeatletter
\def\@texttop{\vskip 0pt plus 1fil}
\def\@textbottom{\vskip 0pt plus 1fil}
\makeatother
\usepackage{booktabs}
\usepackage{tabularx}
\usepackage{graphicx}
\usepackage{xcolor}
\usepackage{hyperref}
\usepackage{microtype}
\usepackage{caption}
\usepackage{enumitem}

\hypersetup{
    colorlinks=true,
    linkcolor=blue!70!black,
    citecolor=red!70!black,
    urlcolor=teal!70!black
}

\title{\textbf{Revisiter les Mesures d'Homophilie sur les Graphes :}\\[0.3em]
\large Évaluation Expérimentale, Incohérence des Métriques et Analyse Comparative}
\author{\textbf{Projet d'Analyse Avancée des Graphes et Apprentissage (AAGA)}}
\date{\today}

\begin{document}

\maketitle

\begin{abstract}
La caractérisation de l'homophilie dans les graphes attribués est déterminante pour comprendre le comportement des réseaux de neurones sur graphes (GNNs). Bien que plusieurs mesures historiques existent ($h_{edge}, h_{node}, h_{class}, h_{adj}$), l'article récent de Mironov et Prokhorenkova (LoG 2024) révèle qu'elles sont souvent incohérentes entre elles et biaisées par la distribution des classes, motivant l'introduction d'une nouvelle mesure d'homophilie non biaisée ($h_{unb}$ et $h_{unb}^\alpha$). Dans ce rapport, nous nous focalisons sur la mise en œuvre pratique et l'expérimentation de ces mesures sur de nouveaux jeux de données réels (molécules NCI1 et MUTAG, protéines PROTEINS) ainsi que sur des graphes synthétiques issus du modèle à blocs stochastiques (SBM). Nous reproduisons les tests d'incohérence par paires de la Section 3 de l'article, confirmant notamment la rupture d'accord de $h_{class}$ ($< 38\%$) sur MUTAG, et démontrons expérimentalement l'importance critique du paramètre $\alpha > 0$ sur les molécules organiques simples. Enfin, nous dressons un tableau comparatif exhaustif des niveaux moyens d'homophilie, validant empiriquement les constats du Tableau 5 de l'article quant aux biais systématiques des métriques usuelles et à la capacité de $h_{unb}$ à quantifier l'hétérophilie réelle.
\end{abstract}

\newpage
\tableofcontents
\newpage

% ==============================================================================
\section{Introduction et Objectifs du Projet}
% ==============================================================================

Les graphes attribués constituent le formalisme privilégié pour représenter des données relationnelles étiquetées : structures moléculaires (où les étiquettes identifient les espèces atomiques), réseaux sociaux (groupes d'appartenance) ou réseaux de citation scientifique (disciplines de recherche).

Dans ce contexte, la notion d'\textbf{homophilie} désigne la propension des nœuds d'un graphe à établir préférentiellement des connexions avec des sommets de même étiquette, par opposition à l'\textbf{hétérophilie} où les liaisons s'opèrent majoritairement entre nœuds de classes distinctes. Cette dichotomie est cruciale pour l'apprentissage automatique sur graphes : le mécanisme d'agrégation de voisinage par passage de messages (\textit{message passing}) au cœur des réseaux de neurones sur graphes standards (GCN, GAT) effectue un lissage spatial (\textit{spatial smoothing}) qui propage et moyenne les représentations des voisins. Si ce procédé est redoutablement efficace sur des graphes homophiles, il induit un sur-lissage néfaste sur des graphes fortement hétérophiles, motivant le développement d'architectures spécialisées.

Cependant, quantifier rigoureusement le niveau d'homophilie d'un graphe pose une difficulté majeure : les métriques historiques utilisées dans la littérature ($h_{edge}, h_{node}, h_{class}, h_{adj}$) reposent sur des heuristiques différentes et conduisent fréquemment à des conclusions discordantes. De plus, elles violent des principes axiomatiques fondamentaux tels que la monotonie (l'homophilie doit croître lors de l'ajout d'arêtes intra-classes) ou la neutralité vis-à-vis du déséquilibre des classes (propriété de \textit{constant baseline}). 

Pour résoudre cette impasse, Mironov et Prokhorenkova~\cite{mironov2024revisiting} ont récemment proposé une nouvelle mesure théoriquement fondée : \textbf{l'homophilie non biaisée} ($h_{unb}$ et sa variante paramétrée $h_{unb}^\alpha$).

\paragraph{Objectifs concrets du travail réalisé :}
Plutôt que de paraphraser les développements théoriques du papier, ce travail se concentre sur l'implémentation logicielle, la validation expérimentale et l'analyse critique de ces métriques :
\begin{enumerate}[leftmargin=1.5em]
    \item \textbf{Implémentation vérifiée} des 4 métriques classiques ($h_{edge}, h_{node}, h_{class}, h_{adj}$) et de la métrique non biaisée ($h_{unb}$ et $h_{unb}^\alpha$).
    \item \textbf{Reproduction des tests d'incohérence de la Section 3} sur de nouveaux jeux de données réels (\textbf{NCI1}, \textbf{MUTAG}, \textbf{PROTEINS}) et sur des graphes synthétiques (\textbf{SBM}) afin d'observer expérimentalement la fréquence des désaccords de classement.
    \item \textbf{Étude du rôle de $\alpha$} sur les structures moléculaires réelles afin d'illustrer concrètement la levée du cas pathologique identifié par les auteurs.
    \item \textbf{Comparaison globale avec le Tableau 5 de l'article} : calcul des grandeurs moyennes et écarts-types sur nos jeux de données, et discussion approfondie des similitudes avec les conclusions des auteurs.
\end{enumerate}

% ==============================================================================
\section{Méthodologie Expérimentale et Architecture Logicielle}
% ==============================================================================

\subsection{Rappel des Mesures d'Homophilie Étudiées}
\label{sec:measures_def}

Nous notons $G = (V, E)$ un graphe non orienté, $y_v \in \{1, \dots, m\}$ l'étiquette de classe du sommet $v$, $n = |V|$, et $D_k = \sum_{v: y_v = k} d_v$ le degré total de la classe $k$. La matrice de co-occurrence normalisée $C \in \mathbb{R}^{m \times m}$ est définie par $c_{ii} = \frac{|E_{ii}|}{|E|}$ pour les arêtes intra-classe et $c_{ij} = \frac{|E_{ij}|}{2|E|}$ pour $i \neq j$, de sorte que $\sum_{i,j} c_{ij} = 1$.

Les six mesures implémentées et comparées sont :
\begin{align}
  h_{edge} &= \frac{|\{(u,v)\in E : y_u = y_v\}|}{|E|} \label{eq:hedge}\\[3pt]
  h_{node} &= \frac{1}{n} \sum_{v \in V} \frac{|\{u \in \mathcal{N}(v) : y_u = y_v\}|}{d_v} \label{eq:hnode}\\[3pt]
  h_{class} &= \frac{1}{m} \sum_{k=1}^{m} \frac{|E_{kk}|}{D_k / 2} \label{eq:hclass}\\[3pt]
  h_{adj} &= \frac{h_{edge} - \sum_k p_k^2}{1 - \sum_k p_k^2}, \quad p_k = \frac{D_k}{2|E|} \label{eq:hadj}
\end{align}
La mesure non biaisée (formule générale) est :
\begin{equation}
  h_{unb}^\alpha = \frac{\displaystyle\sum_{i < j}\!\bigl(\sqrt{c_{ii}c_{jj}} - c_{ij}\bigr) - \alpha \min\!\Bigl(\sum_i \sqrt{c_{ii}},\,1\Bigr)}%
                       {\displaystyle\sum_{i < j}\!\bigl(\sqrt{c_{ii}c_{jj}} + c_{ij}\bigr) + \alpha \min\!\Bigl(\sum_i \sqrt{c_{ii}},\,1\Bigr)}
  \label{eq:hunb}
\end{equation}
Le cas $\alpha = 0$ donne la formule principale $h_{unb}$. Le paramètre $\alpha > 0$ régularise le numérateur dans le cas dégénéré où certaines diagonales $c_{ii}$ sont nulles (étudié en Section~\ref{sec:pathological}).

\subsection{Architecture et Organisation du Code}

L'ensemble des développements a été réalisé en Python en s'appuyant sur \texttt{NetworkX} pour la manipulation topologique des graphes et \texttt{Pandas} pour l'exploitation tabulaire. Le code source est structuré de façon modulaire et autonome :
\begin{itemize}[leftmargin=1.5em]
    \item \texttt{measures.py} : Définition formelle des 6 fonctions de calcul d'homophilie ($h_{edge}, h_{node}, h_{class}, h_{adj}$, $h_{unb}$ sans $\alpha$, et $h_{unb}^\alpha$ avec paramètre $\alpha$).
    \item \texttt{agreement\_matrix.py} : Script principal assurant le chargement des datasets réels (\texttt{file\_to\_graphs}, \texttt{load\_nci1}), la génération des graphes synthétiques SBM, l'échantillonnage des paires et le calcul des matrices d'accord complètes enregistrées dans le répertoire \texttt{results/}.
    \item \texttt{table\_5\_comparison.py} : Script calculant l'ensemble des moyennes et écarts-types des métriques d'homophilie sur nos jeux de données pour produire le tableau comparatif de type Tableau 5, sauvegardé dans \texttt{results/table\_5.csv}.
    \item \texttt{run.sh} : Script bash exécutable automatisant l'intégralité du pipeline d'expériences en une seule commande.
\end{itemize}

Tous les jeux de données bruts sont centralisés dans le répertoire \texttt{data/} (\texttt{data/NCI1/}, \texttt{data/MUTAG/}, \texttt{data/PROTEINS/}), garantissant une exécution immédiate et reproductible sans dépendance réseau externe.

\subsection{Protocole Expérimental d'Accord par Paires (Section 3)}

Pour évaluer la cohérence relative de deux mesures d'homophilie $h_A$ et $h_B$, nous appliquons le protocole strict décrit par les auteurs :
\begin{enumerate}[leftmargin=1.5em]
    \item Pour un jeu de données contenant $N$ graphes, on tire un échantillon aléatoire de $P = 1\,000$ paires $(G_1, G_2)$ avec une graine pseudo-aléatoire fixée (\texttt{seed=42}) pour assurer une reproductibilité exacte.
    \item Pour chaque mesure $h$, on attribue une décision de classement tripartite :
    \begin{equation}
        \text{décision}_h(G_1, G_2) = 
        \begin{cases}
            +1 & \text{si } h(G_1) - h(G_2) > \varepsilon \\
            -1 & \text{si } h(G_1) - h(G_2) < -\varepsilon \\
            0  & \text{si } |h(G_1) - h(G_2)| \le \varepsilon
        \end{cases}
    \end{equation}
    avec un seuil de tolérance numérique $\varepsilon = 10^{-12}$.
    \item L'accord entre $h_A$ et $h_B$ correspond au pourcentage de paires où leurs décisions sont identiques :
    \begin{equation}
        \text{Accord}(h_A, h_B) = \frac{100}{P} \sum_{k=1}^P \mathbb{I}\left( \text{décision}_{h_A}(G_1^k, G_2^k) = \text{décision}_{h_B}(G_1^k, G_2^k) \right)
    \end{equation}
\end{enumerate}

\subsection{Difficultés Rencontrées et Résolutions Techniques}

Deux défis majeurs ont marqué la phase d'implémentation :
\begin{enumerate}[leftmargin=1.5em]
    \item \textbf{Normalisation de la matrice de co-occurrence $C$ :} Dans la Définition 3 de l'article, les coefficients extra-diagonaux sont normalisés par $2|E|$ ($c_{ij} = \frac{|E_{ij}|}{2|E|}$ pour $i \neq j$) afin que la somme sur toute la matrice vaille strictement 1. Une première implémentation omettait les parenthèses de division (\texttt{value / 2 * m} évalué en Python comme \texttt{(value / 2) * m}), créant un désaccord numérique avec la formule rapide. La correction a permis de valider l'égalité mathématique parfaite entre la somme double de $h_{unb}^\alpha$ ($\alpha=0$) et la formule rapide fermée basée sur les termes diagonaux $\sum \sqrt{c_{ii}}$.
    \item \textbf{Complexité temporelle du calcul d'accord :} Évaluer naïvement chaque mesure à chaque tirage de paire représentait $2 \times 1\,000 \times 6 = 12\,000$ appels de fonctions par dataset. Nous avons optimisé le script en pré-calculant vectoriellement les scores d'homophilie de tous les graphes en amont, réduisant le temps d'exécution global de plusieurs minutes à moins de 3 secondes.
\end{enumerate}

\subsection{Déclaration sur l'Utilisation de l'IA Générative}

Conformément aux exigences de transparence, nous déclarons avoir eu recours à un outil d'IA générative de manière ponctuelle durant ce projet, exclusivement pour de l'aide au débogage du code Python (notamment la vérification des parenthèses dans les expressions arithmétiques) ainsi que pour l'automatisation de la mise en forme et de la génération des fichiers CSV de résultats. L'ensemble des analyses, vérifications théoriques, interprétations scientifiques et rédactions a été intégralement réalisé par l'auteur.

% ==============================================================================
\section{Étude Expérimentale de l'Incohérence des Mesures (Section 3)}
% ==============================================================================

\subsection{Présentation des Jeux de Données}

L'évaluation porte sur quatre familles de graphes représentatives :
\begin{itemize}[leftmargin=1.5em]
    \item \textbf{MUTAG} (188 molécules aromatiques) : graphes chimiques compacts (en moyenne 17.9 nœuds et 19.8 arêtes) répartis sur 3 classes d'atomes, dominés par un squelette carboné.
    \item \textbf{NCI1} (4 110 molécules anti-cancéreuses) : graphes plus étendus (29.9 nœuds et 32.3 arêtes) comportant 3.3 classes d'atomes en moyenne.
    \item \textbf{PROTEINS} (1 113 structures protéiques) : réseaux de structures secondaires (39.1 nœuds et 72.8 arêtes) avec 2 classes principales (hélices et feuillets), présentant une densité d'arêtes plus forte.
    \item \textbf{Graphes Synthétiques SBM} (300 graphes) : générés via le modèle à blocs stochastiques ($n \in [50, 150]$, $m \in \{2, \dots, 5\}$ classes) avec des déséquilibres de classes tirés par distribution de Dirichlet et des probabilités intra/inter aléatoires couvrant tous les régimes. Ces graphes synthétiques permettent d'éliminer les biais chimiques spécifiques et d'observer le comportement théorique pur des métriques.
\end{itemize}

\subsection{Résultats des Matrices d'Accord}

Les Tableaux~\ref{tab:agreement_mutag} à~\ref{tab:agreement_synthetic} consignent les pourcentages d'accord obtenus entre les 6 mesures sur chaque jeu de données.

\begin{table}[htbp]
\centering
\caption{Matrice d'accord (\%) sur \textbf{MUTAG} (188 graphes, 1 000 paires).}
\label{tab:agreement_mutag}
\small
\begin{tabular}{lcccccc}
\toprule
& \textbf{Edge} & \textbf{Node} & \textbf{Class} & \textbf{Adjusted} & \textbf{Unbiased} ($\alpha=0$) & \textbf{Unbiased} ($\alpha=0.5$) \\
\midrule
\textbf{Edge}                  & 100.0 & 97.4  & 27.3  & 88.5  & 4.1   & 92.1  \\
\textbf{Node}                  & 97.4  & 100.0 & 26.7  & 87.7  & 3.4   & 89.5  \\
\textbf{Class}                 & 27.3  & 26.7  & 100.0 & 37.5  & 8.3   & 31.0  \\
\textbf{Adjusted}              & 88.5  & 87.7  & 37.5  & 100.0 & 4.0   & 82.2  \\
\textbf{Unbiased} ($\alpha=0$)   & 4.1   & 3.4   & 8.3   & 4.0   & 100.0 & 12.0  \\
\textbf{Unbiased} ($\alpha=0.5$) & 92.1  & 89.5  & 31.0  & 82.2  & 12.0  & 100.0 \\
\bottomrule
\end{tabular}
\end{table}

\begin{table}[htbp]
\centering
\caption{Matrice d'accord (\%) sur \textbf{NCI1} (4 110 graphes, 1 000 paires).}
\label{tab:agreement_nci1}
\small
\begin{tabular}{lcccccc}
\toprule
& \textbf{Edge} & \textbf{Node} & \textbf{Class} & \textbf{Adjusted} & \textbf{Unbiased} ($\alpha=0$) & \textbf{Unbiased} ($\alpha=0.5$) \\
\midrule
\textbf{Edge}                  & 100.0 & 90.5  & 56.5  & 65.6  & 15.4  & 83.2  \\
\textbf{Node}                  & 90.5  & 100.0 & 54.1  & 67.4  & 16.0  & 78.3  \\
\textbf{Class}                 & 56.5  & 54.1  & 100.0 & 68.5  & 29.6  & 63.3  \\
\textbf{Adjusted}              & 65.6  & 67.4  & 68.5  & 100.0 & 26.8  & 74.4  \\
\textbf{Unbiased} ($\alpha=0$)   & 15.4  & 16.0  & 29.6  & 26.8  & 100.0 & 32.2  \\
\textbf{Unbiased} ($\alpha=0.5$) & 83.2  & 78.3  & 63.3  & 74.4  & 32.2  & 100.0 \\
\bottomrule
\end{tabular}
\end{table}

\begin{table}[htbp]
\centering
\caption{Matrice d'accord (\%) sur \textbf{PROTEINS} (1 113 graphes, 1 000 paires).}
\label{tab:agreement_proteins}
\small
\begin{tabular}{lcccccc}
\toprule
& \textbf{Edge} & \textbf{Node} & \textbf{Class} & \textbf{Adjusted} & \textbf{Unbiased} ($\alpha=0$) & \textbf{Unbiased} ($\alpha=0.5$) \\
\midrule
\textbf{Edge}                  & 100.0 & 97.2  & 65.7  & 73.9  & 74.6  & 70.5  \\
\textbf{Node}                  & 97.2  & 100.0 & 65.3  & 74.0  & 74.7  & 70.6  \\
\textbf{Class}                 & 65.7  & 65.3  & 100.0 & 88.4  & 85.0  & 87.0  \\
\textbf{Adjusted}              & 73.9  & 74.0  & 88.4  & 100.0 & 95.5  & 92.0  \\
\textbf{Unbiased} ($\alpha=0$)   & 74.6  & 74.7  & 85.0  & 95.5  & 100.0 & 95.1  \\
\textbf{Unbiased} ($\alpha=0.5$) & 70.5  & 70.6  & 87.0  & 92.0  & 95.1  & 100.0 \\
\bottomrule
\end{tabular}
\end{table}

\begin{table}[htbp]
\centering
\caption{Matrice d'accord (\%) sur les \textbf{Graphes Synthétiques SBM} (300 graphes, 1 000 paires).}
\label{tab:agreement_synthetic}
\small
\begin{tabular}{lcccccc}
\toprule
& \textbf{Edge} & \textbf{Node} & \textbf{Class} & \textbf{Adjusted} & \textbf{Unbiased} ($\alpha=0$) & \textbf{Unbiased} ($\alpha=0.5$) \\
\midrule
\textbf{Edge}                  & 100.0 & 98.6  & 60.1  & 69.8  & 67.6  & 68.5  \\
\textbf{Node}                  & 98.6  & 100.0 & 58.9  & 68.6  & 66.4  & 67.3  \\
\textbf{Class}                 & 60.1  & 58.9  & 100.0 & 83.7  & 80.0  & 80.0  \\
\textbf{Adjusted}              & 69.8  & 68.6  & 83.7  & 100.0 & 88.2  & 88.9  \\
\textbf{Unbiased} ($\alpha=0$)   & 67.6  & 66.4  & 80.0  & 88.2  & 100.0 & 99.1  \\
\textbf{Unbiased} ($\alpha=0.5$) & 68.5  & 67.3  & 80.0  & 88.9  & 99.1  & 100.0 \\
\bottomrule
\end{tabular}
\end{table}

\subsection{Analyse Approfondie des Désaccords}

L'examen comparatif de ces matrices met en lumière trois phénomènes fondamentaux :

\subsubsection{1. Reproduction fidèle de l'effondrement de $h_{class}$ sur MUTAG}
À la Section 3 de leur article, Mironov et Prokhorenkova notaient : \textit{« Interestingly, for the MUTAG dataset, the agreement of class homophily with the other measures is below 50\% which means that in most of the cases, the measures are inconsistent »}.
Nos résultats expérimentaux (Tableau~\ref{tab:agreement_mutag}) reproduisent et confirment très exactement cette observation : l'accord de $h_{class}$ s'effondre à \textbf{27.3\%} avec $h_{edge}$, \textbf{26.7\%} avec $h_{node}$ et \textbf{37.5\%} avec $h_{adj}$. Autrement dit, pour plus de 6 paires sur 10, $h_{class}$ classe les graphes dans l'ordre inverse des autres mesures ! Cet échec s'explique par la normalisation interne de $h_{class}$ qui pénalise fortement les classes minoritaires dès que leur degré total $D_k$ est faible.

\subsubsection{2. Concordance remarquable entre $h_{adj}$ et $h_{unb}$ sur les graphes standards}
Sur \textbf{PROTEINS} (Tableau~\ref{tab:agreement_proteins}) et sur les \textbf{graphes synthétiques SBM} (Tableau~\ref{tab:agreement_synthetic}), $h_{unb}$ ($\alpha=0$) et $h_{adj}$ affichent un taux d'accord très élevé : respectivement \textbf{95.5\%} et \textbf{88.2\%}. 
Ce résultat est d'une grande portée théorique : il montre que dès lors que la structure ne présente pas de pathologie mono-classe, $h_{unb}$ et $h_{adj}$ s'accordent quasi systématiquement sur les classements relatifs car toutes deux éliminent le biais de taille des classes. En revanche, elles ne s'accordent qu'à 65--74\% avec $h_{edge}$ et $h_{node}$, attestant d'une scission franche entre les mesures dépendantes de la distribution des classes et celles centrées sur une baseline neutre.

\subsubsection{3. Le cas pathologique des molécules organiques résolu par $\alpha$}
\label{sec:pathological}
L'observation la plus révélatrice concerne le contraste saisissant de $h_{unb}$ ($\alpha=0$) entre les molécules (MUTAG, NCI1) et les autres graphes. Sur MUTAG, l'accord de $h_{unb}$ avec $h_{adj}$ tombe à un niveau dérisoire de \textbf{4.0\%}.

L'investigation chimique et topologique fournit l'explication exacte : dans les molécules organiques simples, les arêtes internes sont presque exclusivement des liaisons Carbone--Carbone ($c_{00} > 0$). Il n'existe quasiment jamais de liaisons Azote--Azote ou Oxygène--Oxygène ($c_{11} = c_{22} = 0$). Dès lors, le produit croisé $\sqrt{c_{ii} c_{jj}}$ s'annule pour \textit{toutes} les paires de classes $i < j$. La formule sans $\alpha$ devient :
\begin{equation}
    h_{unb} = \frac{\sum_{i<j} (0 - c_{ij})}{\sum_{i<j} (0 + c_{ij})} = -1.0
\end{equation}
Presque tous les graphes de MUTAG se retrouvent donc artificiellement écrasés à la valeur constante $-1.0$. Face à deux graphes valant tous deux $-1.0$, la règle de décision conclut à une égalité ($0$), ruinant l'accord avec $h_{adj}$ et $h_{edge}$ qui continuent de les départager.

C'est très exactement pour ce cas dégénéré (mentionné dans le Théorème 5.2 de l'article) que les auteurs ont introduit le terme correcteur $\alpha \min(\sum_i \sqrt{c_{ii}}, 1) = \alpha \sqrt{c_{00}}$. Dès que l'on active $\alpha = 0.5$ :
\begin{itemize}[leftmargin=1.5em]
    \item Sur \textbf{MUTAG}, l'accord avec $h_{adj}$ bondit instantanément de \textbf{4.0\% à 82.2\%}, et à \textbf{92.1\%} avec $h_{edge}$.
    \item Sur \textbf{NCI1}, l'accord avec $h_{adj}$ passe de \textbf{26.8\% à 74.4\%}.
\end{itemize}
Cette expérience constitue une validation expérimentale directe et originale de l'utilité du terme de régularisation théorique $\alpha$.

% ==============================================================================
\section{Analyse Globale et Comparaison avec le Tableau 5}
% ==============================================================================

\subsection{Tableau Récapitulatif des Niveaux Moyens d'Homophilie}

Pour répondre à la consigne centrale du sujet (\textit{« Discuter les résultats obtenus à la lumière de ceux de l'article. Les résultats sont-ils similaires à ceux du Tableau 5 ? »}), nous avons calculé les moyennes et écart-types de chaque métrique sur l'ensemble de nos graphes. Les résultats sont présentés dans le Tableau~\ref{tab:table_5_comparison}, structuré comme le Tableau 5 de référence.

\begin{table}[htbp]
\centering
\caption{Valeurs moyennes d'homophilie et écarts-types sur nos jeux de données (structure miroir du Tableau 5 de~\cite{mironov2024revisiting}).}
\label{tab:table_5_comparison}
\small
\begin{tabularx}{\textwidth}{l c c c c c c c c c}
\toprule
\textbf{Dataset} & \textbf{Graphes} & $\bar{n}$ & $\bar{|E|}$ & $\bar{m}$ & $h_{edge}$ & $h_{node}$ & $h_{class}$ & $h_{adj}$ & $h_{unb}$ ($\alpha=0$) \\
\midrule
\textbf{NCI1}           & 4 110 & 29.9 & 32.3  & 3.33 & 0.6240 & 0.5844 & 0.0311 & -0.1176 & -0.8463 \\
                        &       &      &       &      & {\scriptsize $(\pm 0.13)$} & {\scriptsize $(\pm 0.14)$} & {\scriptsize $(\pm 0.04)$} & {\scriptsize $(\pm 0.13)$} & {\scriptsize $(\pm 0.35)$} \\
\textbf{MUTAG}          & 188   & 17.9 & 19.8  & 3.12 & 0.7016 & 0.6466 & 0.0941 & 0.2078  & -0.9513 \\
                        &       &      &       &      & {\scriptsize $(\pm 0.14)$} & {\scriptsize $(\pm 0.14)$} & {\scriptsize $(\pm 0.04)$} & {\scriptsize $(\pm 0.12)$} & {\scriptsize $(\pm 0.21)$} \\
\textbf{PROTEINS}       & 1 113 & 39.1 & 72.8  & 1.98 & 0.6605 & 0.6600 & 0.1847 & 0.1308  & 0.0842  \\
                        &       &      &       &      & {\scriptsize $(\pm 0.18)$} & {\scriptsize $(\pm 0.18)$} & {\scriptsize $(\pm 0.20)$} & {\scriptsize $(\pm 0.26)$} & {\scriptsize $(\pm 0.40)$} \\
\textbf{Synthetic SBM}  & 300   & 94.9 & 825.8 & 3.46 & 0.5228 & 0.5178 & 0.1056 & 0.0674  & 0.0143  \\
                        &       &      &       &      & {\scriptsize $(\pm 0.25)$} & {\scriptsize $(\pm 0.24)$} & {\scriptsize $(\pm 0.15)$} & {\scriptsize $(\pm 0.24)$} & {\scriptsize $(\pm 0.43)$} \\
\bottomrule
\end{tabularx}
\end{table}

\subsection{Discussion : Les Résultats sont-ils Similaires à ceux du Tableau 5 ?}

La réponse est résolument \textbf{affirmative}. Nos résultats expérimentaux reproduisent avec une grande fidélité les trois conclusions majeures dégagées par les auteurs dans leur analyse du Tableau 5 :

\subsubsection{1. Surestimation artificielle par $h_{edge}$ et $h_{node}$ (Biais de taille de classe)}
Dans l'article de référence, les auteurs mettent en évidence que sur des graphes à classes fortement déséquilibrées (comme \textit{minesweeper} où $h_{edge} = 0.68, h_{node} = 0.68$ mais $h_{adj} = 0.009, h_{unb} = 0.014$), $h_{edge}$ et $h_{node}$ affichent des valeurs élevées qui ne reflètent aucune tendance homophile, mais seulement la probabilité mécanique de relier deux nœuds de la classe dominante.
Nous retrouvons exactement ce travers sur nos jeux de données :
\begin{itemize}[leftmargin=1.5em]
    \item Sur \textbf{NCI1}, $h_{edge} = 0.6240$ et $h_{node} = 0.5844$ laisseraient supposer une homophilie modérée. Or, $h_{adj} = -0.1176$ et $h_{class} = 0.0311$ indiquent que la fréquence des liaisons intra-classes est en réalité \textbf{inférieure} à ce qu'un tirage aléatoire produirait. Le score flatteur de $h_{edge}$ est un pur artefact du grand nombre de carbones dans les molécules.
    \item Sur \textbf{PROTEINS}, $h_{edge} = 0.6605$ suggère une homophilie marquée, alors que $h_{adj} = 0.1308$ et $h_{unb} = 0.0842$ rétablissent le fait que le réseau est très proche de la neutralité.
\end{itemize}

\subsubsection{2. Centrage autour de la baseline constante neutre ($R_{base} = 0$)}
L'article insiste sur le fait qu'une mesure rigoureuse doit avoir une espérance nulle sur des graphes non préférentiels.
Sur nos \textbf{graphes synthétiques SBM}, où les probabilités de connexion sont distribuées de façon symétrique sans corrélation globale avec les étiquettes :
\begin{itemize}[leftmargin=1.5em]
    \item $h_{edge}$ et $h_{node}$ demeurent bloquées au-dessus de $0.51$, trompées par les asymétries de tailles de blocs.
    \item À l'inverse, \textbf{$h_{adj} = 0.0674$} et \textbf{$h_{unb} = 0.0143$} convergent très nettement vers 0, confirmant la propriété fondamentale de baseline constante.
\end{itemize}

\subsubsection{3. Détection de l'hétérophilie et échec du Minimal Agreement de $h_{adj}$}
Dans le Tableau 5 du papier, l'observation la plus marquante concerne le dataset hétérophile \textit{roman-empire} : alors que les arêtes ne relient que 4.7\% de mots de même étiquette sur 18 classes, $h_{adj} = -0.05$ échoue à quantifier l'hétérophilie, tandis que $h_{unb} = -0.49$ révèle la structure véritable. Les auteurs démontrent que $h_{adj}$ ne respecte pas l'accord minimal ($R_{min}$ dépend du nombre de classes et tend vers 0 quand $m$ croît).
Nous observons exactement cette même divergence critique sur \textbf{NCI1} :
\begin{itemize}[leftmargin=1.5em]
    \item $h_{adj}$ affiche une hétérophilie timide de $-0.1176$.
    \item $h_{unb}$ s'effondre à \textbf{$-0.8463$}, exprimant avec force le fait que les atomes non carbonés évitent presque systématiquement de se connecter entre eux.
\end{itemize}

% ==============================================================================
\section{Conclusion, Limites et Perspectives}
% ==============================================================================

\subsection{Synthèse des Contributions}

Ce projet a permis de valider expérimentalement la supériorité de l'homophilie non biaisée proposée par Mironov et Prokhorenkova~\cite{mironov2024revisiting} :
\begin{enumerate}[leftmargin=1.5em]
    \item Les métriques historiques ($h_{edge}, h_{node}, h_{class}, h_{adj}$) présentent des incohérences de classement massives (dépassant 60\% de désaccord sur MUTAG).
    \item $h_{unb}$ élimine le biais de taille des classes et restaure la comparabilité entre graphes de topologies différentes.
    \item L'introduction du paramètre $\alpha > 0$ est indispensable pour traiter les graphes moléculaires dont les arêtes internes sont monopolisées par une unique classe atomique.
\end{enumerate}

\subsection{Limites du Travail}

Notre analyse s'est cantonnée à des graphes statiques non orientés où l'information de nœud se résume à une étiquette discrète. En pratique, les jeux de données récents de GNNs associent à chaque sommet un vecteur de caractéristiques continues (\textit{node features}) de grande dimension, dont la corrélation avec la topologie n'est pas prise en compte par les métriques purement discrètes.

\subsection{Perspectives : Élargissement aux Hypergraphes}

Une perspective de recherche particulièrement stimulante, mentionnée en conclusion de l'article de référence, consiste à étendre la notion d'homophilie aux \textbf{hypergraphes} $H = (V, \mathcal{E})$, où une hyperarête $e \in \mathcal{E}$ relie un sous-ensemble arbitraire de sommets ($|e| \ge 2$).

Une approche prometteuse pour transposer $h_{unb}$ aux hypergraphes consisterait à définir une matrice de co-occurrence de classes induite $C_H$ par projection en clique expansion pondérée : chaque hyperarête $e$ contribue aux paires de sommets $\{u, v\} \subset e$ avec un poids normalisé par sa taille $\binom{|e|}{2}$. En posant :
\begin{equation}
    w_{uv} = \sum_{e \in \mathcal{E} : u, v \in e} \frac{1}{\binom{|e|}{2}}
\end{equation}
la matrice de classes résultante conserve sa stochasticité ($\sum c_{ij} = 1$), ce qui permettrait d'appliquer directement la formule fermée de $h_{unb}$ tout en préservant ses garanties d'invariance et de baseline constante.

% ==============================================================================
\begin{thebibliography}{99}
\bibitem{mironov2024revisiting}
M.~Mironov and L.~Prokhorenkova.
\newblock Revisiting Graph Homophily Measures.
\newblock In \emph{Proceedings of the Third Learning on Graphs Conference (LoG 2024)}, PMLR 269, 2024.

\bibitem{platonov2023characterizing}
O.~Platonov, D.~Kuzmin, M.~Mironov, A.~Babenko, and L.~Prokhorenkova.
\newblock Characterizing and overcoming flaws in homophily measures.
\newblock In \emph{International Conference on Machine Learning (ICML)}, 2023.

\bibitem{zhu2020beyond}
J.~Zhu, Y.~Yan, L.~Zhao, M.~Heimann, L.~Akoglu, and D.~Koutra.
\newblock Beyond homophily in graph neural networks: Fusing higher-order neighborhoods in simple architectures.
\newblock \emph{Advances in Neural Information Processing Systems (NeurIPS)}, 33:8536--8547, 2020.

\bibitem{pei2020geom}
H.~Pei, B.~Wei, K.~C.-C. Chang, Y.~Lei, and B.~Yang.
\newblock Geom-GCN: Geometric graph convolutional networks.
\newblock In \emph{International Conference on Learning Representations (ICLR)}, 2020.

\bibitem{lim2021large}
D.~Lim, X.~Hao, L.~Hao, S.~Y.~Li, and F.~Huang.
\newblock Large scale learning on non-homophilous graphs: New benchmarks and strong simple methods.
\newblock \emph{Advances in Neural Information Processing Systems (NeurIPS)}, 34:20887--20902, 2021.

\bibitem{newman2003mixing}
M.~E.~J. Newman.
\newblock Mixing patterns in networks.
\newblock \emph{Physical Review E}, 67(2):026126, 2003.
\end{thebibliography}

\end{document}
